\documentclass[11pt]{report}
\input{../pri-end/T0_preamble_standalone_De}
\usepackage{longtable}
\usepackage{booktabs}
\hypersetup{
	pdftitle={254 · FFGFT — Zwei Ordnungsprinzipien und das Flächen\-schranken-Theorem},
	pdfauthor={Johann Pascher},
	pdfsubject={Resonanzprinzip, Entropie, Holographische Schranke, Vopson, Infodynamik, T4-Geometrie}
}
\fancyhead[L]{\small\textcolor{t0gray}{FFGFT · Zwei Ordnungsprinzipien}}
\fancyhead[R]{\small\textcolor{t0gray}{Johann Pascher · 2026}}
\title{%
	{\LARGE\bfseries\color{t0blue}
	 Zwei Ordnungsprinzipien und das Flächen\-schranken-Theorem}\\[0.5em]
	{\large Resonanz, Entropie und ihre geometrische Auflösung}\\[1em]
	{\normalsize Johann Pascher \quad $\cdot$ \quad Mai 2026}\\[0.5em]
	{\small \url{https://github.com/jpascher/T0-Time-Mass-Duality}}\\[0.3em]
	{\small Zenodo: \href{https://doi.org/10.5281/zenodo.20117635}{10.5281/zenodo.20117635} (v1.1.0)}
}
\date{Mai 2026}
\begin{document}
	\maketitle
	% ==============================================================================
% 254_Zwei_Ordnungsprinzipien_De_ch.tex
% Kapitelinhalt – einzubinden via Wrapper
% Autor: Johann Pascher · Mai 2026
% ==============================================================================

\section*{Abstract}

In der Physik schwarzer Löcher treten zwei Ordnungsprinzipien auf, die sich scheinbar widersprechen: das \textbf{Resonanzprinzip} (die Natur selektiert resonante Konfigurationen, ξ-Rekursion) und das \textbf{thermodynamische Prinzip} (Entropie strebt nach Maximum, Ausgleich). Dieses Dokument zeigt, dass der scheinbare Widerspruch durch eine Schichtentrennung aufgelöst wird --- die zwei Prinzipien operieren auf verschiedenen Ebenen --- und dass sie beim Verdampfungs-Endpunkt nicht nur verträglich, sondern gegenseitig erzwingend sind. Ein zentrales geometrisches Theorem (Flächen\-schranken-Argument) wird hergeleitet. Die Verbindung zu Vopsons zweitem Infodynamik-Gesetz wird als ehrliche Übersetzung --- gleiche Phänomenologie, verschiedene Mechanismen --- ausgewiesen. Drei offene Beweislücken werden explizit benannt.

\tableofcontents
\newpage

% ------------------------------------------------------------------------------
\section{Die zwei Ordnungsprinzipien}

\subsection{Das Resonanzprinzip}

In FFGFT selektiert die ξ-Rekursion auf dem 4D-Identifikationstorus $T^4$ die zulässigen Windungskonfigurationen. Nur Konfigurationen die die Rekursionsbedingung erfüllen, sind stabil --- Teilchen haben diskrete Massen, keine Kontinua. Das ξ-Prinzip ist ordnungsschaffend: es erzeugt Struktur, bevorzugt den resonanten Weg, und definiert den zulässigen Zustandsraum bevor Dynamik operiert (vgl.\ Dok.~201, Dok.~250).

\subsection{Das thermodynamische Prinzip}

Die Thermodynamik strebt nach maximaler Entropie: Unordnung nimmt zu, Strukturen gleichen sich aus, alles fließt zum gemittelten Zustand. Das erscheint dem Resonanzprinzip direkt entgegengesetzt.

\begin{wichtig}[Der scheinbare Widerspruch]
Das Resonanzprinzip bevorzugt Ordnung und Struktur. Das thermodynamische Prinzip bevorzugt Ausgleich und Unordnung. Beide scheinen gleichzeitig in der Natur zu wirken. Wie ist das möglich?
\end{wichtig}

% ------------------------------------------------------------------------------
\section{Das Flächen\-schranken-Theorem}

\subsection{Die Argumentkette}

Das folgende Argument ist ein \textbf{Theorem}, kein Postulat. Es folgt aus drei etablierten Größen:

\begin{enumerate}
  \item \textbf{Holographische Schranke (Bekenstein):} Die maximale Information die eine Fläche tragen kann, ist proportional zur Fläche: $N_\text{max} \propto A$.
  \item \textbf{Hawking-Verdampfung:} $A \propto M^2$ schrumpft mit der Masse.
  \item \textbf{FFGFT-Untergrenze:} $L_0 = \xipar \cdot \ell_P$ setzt eine minimale Länge, also eine minimale Fläche --- der Zustandsraum endet bei $N = 0$.
\end{enumerate}

\begin{keyresult}[Flächen\-schranken-Theorem]
Aus (1), (2) und (3) folgt zwingend: Die \textbf{mögliche Informationskapazität} der Horizontfläche muss bei Verdampfung sinken. Das ist nicht postuliert, sondern folgt aus der Geometrie allein.
\begin{equation}
  N_\text{max}(M) \propto M^2 \quad \Rightarrow \quad \frac{dN_\text{max}}{dt} < 0 \quad \text{(solange BH verdampft)}
\end{equation}
\end{keyresult}

\subsection{Die Verdichtung als Konsequenz}

Aus dem Theorem folgt, wenn man Erhaltung der topologischen Ladung $Q$ hinzufügt:

\begin{erkenntnis}[Erzwungene Verdichtung]
\begin{itemize}
  \item Die \emph{mögliche} Information (Kapazität $N_\text{max}$) sinkt --- geometrisch erzwungen.
  \item Die \emph{tatsächliche} topologische Ladung $Q$ bleibt erhalten --- topologische Invarianz.
  \item Also muss $Q$ in immer weniger Kapazität gepackt werden: \textbf{Verdichtung ist erzwungen}.
\end{itemize}
Am Remnant ($N = 0$): maximale Ladungskonzentration bei minimaler Kapazität. Das ist das Resonanzprinzip --- nicht gewählt, sondern unvermeidlich.
\end{erkenntnis}

\subsection{Entropie und Resonanz als zwei Seiten eines Prozesses}

Das thermodynamische Prinzip treibt die Verdampfung (Strahlungsentropie steigt, $dS_\text{rad}/dt > 0$). Dabei schrumpft die Fläche. Das Flächenschranken-Theorem zwingt dann die Verdichtung. Entropie und Resonanzprinzip sind also nicht Gegner:

\begin{erkenntnis}[Keine Widersprüchlichkeit, sondern Kausalkette]
Das entropische Prinzip erzeugt die Bedingungen unter denen das Resonanzprinzip unvermeidlich wird. Das eine verursacht das andere.
\end{erkenntnis}

% ------------------------------------------------------------------------------
\section{Python-verifizierte Schlüsselzahlen}

\subsection{Gegenläufigkeit N\_Knoten vs.\ Strahlungsentropie}

Numerische Verifikation (Skript \texttt{infodynamik.py}):

\begin{equation}
  N_\text{Knoten} \propto M^2 \quad \Rightarrow \quad \frac{dN_\text{Knoten}}{dt} < 0 \quad \text{immer}
\end{equation}
\begin{equation}
  \frac{dS_\text{rad}}{dt} > 0 \quad \text{(Page-Faktor 4/3: Strahlung trägt mehr Entropie weg als BH verliert)}
\end{equation}

Die Gegenläufigkeit ist exakt und rechenbar. Für $M = M_\odot$:
\begin{equation}
  \frac{k_B T_H}{E_\text{max}} \approx 5{,}8 \times 10^{-44}, \qquad N_\text{Knoten}(M_\odot) \approx 1{,}5 \times 10^{77}
\end{equation}

Am Remnant ($M_\text{krit} \approx 1{,}15 \times 10^{-13}$ kg): $N_\text{Knoten} \approx 5 \times 10^{-10}$ --- die Kapazität ist praktisch null, die Ladung maximal verdichtet.

\subsection{Skalierungsvergleich: Kapazität vs.\ naive Ladungszählung}

Testet man die Konsistenzbedingung $Q \leq N_\text{max}$ mit der naiven Zählung $Q = M/m_e$ (Skript \texttt{kapazitaet.py}):

\begin{center}
\begin{tabular}{lll}
\toprule
$M$ & $N_\text{max}/Q$ & Status \\
\midrule
$1\,M_\odot$ & $6{,}9 \times 10^{16}$ & konsistent \\
$10^{20}$ kg & $3{,}5 \times 10^{6}$ & konsistent \\
$10^{11}$ kg & $3{,}5 \times 10^{-3}$ & Schere überkreuzt \\
\bottomrule
\end{tabular}
\end{center}

\begin{wichtig}[Bedeutung des Überkreuzens]
Die naive Zählung $Q = M/m_e$ überschreitet $N_\text{max}$ bei $M \approx 2{,}9 \times 10^{13}$ kg --- das ist \emph{nicht} ein physikalischer Widerspruch, sondern der Beweis, dass $M/m_e$ das \textbf{falsche Konto} ist. Es zählt eingefallene Teilchen, nicht topologische Invarianten. Die Konten sind nicht dasselbe.
\end{wichtig}

Daraus folgt eine prüfbare Bedingung: Die korrekte topologische Ladung $Q$ muss höchstens so schnell wachsen wie $M^2$ (die Kapazität), damit $Q \leq N_\text{max}$ immer gilt. Das ist eine Einschränkung an FFGFT, keine freie Wahl.

% ------------------------------------------------------------------------------
\section{Verbindung zu Vopsons Infodynamik}

\subsection{Vopsons zweites Gesetz}

Vopsons zweites Infodynamik-Gesetz besagt: Informationsentropie eines Systems nimmt ab oder bleibt konstant, während thermodynamische Entropie steigt. Die zwei Entropien laufen entgegengesetzt.

\subsection{FFGFT reproduziert die Phänomenologie}

In FFGFT-Größen:
\begin{itemize}
  \item Sinkende Größe: $N_\text{max}(M) \propto M^2$ --- die Kapazität fällt.
  \item Steigende Größe: $S_\text{rad}$ --- thermodynamische Entropie der Strahlung steigt.
\end{itemize}

Die Richtungen sind dieselben wie bei Vopsons Gesetz. \textbf{Aber der Mechanismus ist ein anderer}:

\begin{center}
\begin{tabular}{p{5cm}p{8cm}}
\toprule
\textbf{Vopson} & \textbf{FFGFT} \\
\midrule
Informationsentropie sinkt als \emph{postuliertes Naturgesetz} & Kapazität $N_\text{max}$ sinkt als \emph{geometrisches Theorem} (Fläche + Bekenstein + $L_0$) \\
Mechanismus: Information hat Masse, treibt Minierung & Mechanismus: Flächenschranke + Erhaltung + Untergrenze \\
Zählt: Informations-tragende Teilchenzustände & Zählt: Oberflächenknoten $\propto A \propto M^2$ \\
\bottomrule
\end{tabular}
\end{center}

\begin{keyresult}[Ehrliche Übersetzung]
FFGFT reproduziert die \emph{Phänomenologie} von Vopsons Gesetz über einen anderen Mechanismus. Die zwei Konten zählen Verschiedenes und konvergieren nur im Verhalten. Die Verbindung ist eine \textbf{Übersetzung}, keine Identität.

Dies ist eine stärkere epistemische Position als Vopsons: FFGFT postuliert nichts, sondern leitet die Gegenläufigkeit aus der Geometrie her.
\end{keyresult}

% ------------------------------------------------------------------------------
\section{Drei offene Beweislücken}

Diese werden explizit benannt, weil claim hygiene die Pflicht zur Selbstbegrenzung ist.

\subsection{Beweislücke 1: Die richtige topologische Ladung}

\textbf{Offene Frage:} Was ist die erhaltene Größe $Q$ in Einheiten die mit $N_\text{max}$ vergleichbar sind?

$Q = M/m_e$ ist nachweislich falsch (überschreitet die Kapazität). $Q$ muss höchstens $\propto M^2$ skalieren. Ob $Q$ mit der Bekenstein-Entropie selbst identisch ist (holographisches Prinzip in stärkster Form), oder sub-extensiv, oder etwas anderes --- ist die offene Frage.

\textbf{Warum das wichtig ist:} Erst wenn $Q$ präzise definiert ist, kann die Konsistenzbedingung $Q \leq N_\text{max}$ über die gesamte Verdampfung als Theorem verifiziert werden --- nicht nur numerisch plausibilisiert.

\subsection{Beweislücke 2: Das Extremalprinzip}

\textbf{Offene Frage:} Ist das Resonanzprinzip als Variationsaussage formulierbar?

Die Fibonacci-Konvergenz $\varphi^{-2n}/\sqrt{5} \approx \xipar$ bei $n \approx 8$ und die Lagrange-Steifigkeit $-\frac{1}{2}m_T^2(\partial m)^2$ deuten auf ein Extremalprinzip hin. Aber der explizite Euler-Lagrange-Satz --- ein Beweis, dass der ξ-Fixpunkt ein Minimum einer Informationsfunktion ist --- existiert im Korpus noch nicht.

\textbf{Was der Beweis bräuchte:} Eine Informationsfunktion $\mathcal{I}[\Psi]$ auf dem Konfigurationsraum von $T^4$, so dass $\delta\mathcal{I}/\delta\Psi = 0$ genau die ξ-stabilen Windungskonfigurationen als Lösungen hat.

\subsection{Beweislücke 3: Definition der Informationsentropie in FFGFT}

\textbf{Offene Frage:} Wie definiert FFGFT Informationsentropie präzise?

Bisher: $N_\text{Knoten}$ als Informationsmaß. Aber $N_\text{Knoten}$ ist die holographische Kapazität (Oberflächenknoten), nicht die tatsächlich gespeicherte Information. Die saubere Definition --- was unterscheidet Informationskapazität von Informationsentropie in FFGFT --- fehlt noch. Das ist der Punkt wo eine präzise Brücke zu Vopson möglich wäre, aber noch nicht vollständig gebaut ist.

% ------------------------------------------------------------------------------
\section*{Zusammenfassung}

\begin{longtable}{>{\bfseries}p{6cm} p{8cm}}
\toprule
Punkt & Aussage \\
\midrule
Scheinbarer Widerspruch & Resonanzprinzip (Ordnung) vs.\ thermodynamisches Prinzip (Ausgleich) \\
Auflösung & Verschiedene Ebenen (topologisch vs.\ thermodynamisch) \\
Theorem & Flächenschranke + Bekenstein + $L_0$ $\Rightarrow$ Kapazität sinkt zwingend \\
Kausalkette & Entropie treibt Verdampfung $\Rightarrow$ Fläche schrumpft $\Rightarrow$ Verdichtung erzwungen \\
Vopson-Brücke & Gleiche Phänomenologie, anderer Mechanismus: Übersetzung, nicht Identität \\
Konto-Unterschied & $Q = M/m_e$ (Teilchen) $\neq$ $N_\text{max}$ (Fläche) --- zwei verschiedene Zählungen \\
Offene Frage 1 & Korrekte topologische Ladung $Q$: muss $\leq N_\text{max} \propto M^2$ \\
Offene Frage 2 & Extremalprinzip: ξ-Fixpunkt als Minimum einer Informationsfunktion \\
Offene Frage 3 & Saubere Definition von Informationsentropie in FFGFT \\
\bottomrule
\end{longtable}

	\bibliographystyle{plainnat}
	\bibliography{../bib/FFGFT_Epistemologie}
\end{document}
